\uuid{e80i}
\titre{Raisonnements}
\auteur{bodin ; croizat ; sacre}
\organisation{exo7}
\licence{CC BY-NC-SA 4.0}
\datecreate{2022-03-08}
\niveau{L1}
\module{Algèbre}
\chapitre{Logique, ensemble, raisonnement}
\sousChapitre{Absurde et contraposée}
\difficulte{2}
\theme{Raisonnements}
\exo7qcmid{hilisit.logique.046}

\contenu{
\begin{question}
\qclassification{100.03}
Pour montrer que $\sqrt{3}$ est un nombre irrationnel, je commence une démonstration par l'absurde en écrivant :
\begin{answers}
	\good{Je suppose $\sqrt3\in\Qq$ et je cherche une contradiction.}
	\bad{Je suppose $\sqrt3\notin\Qq$ et je cherche une contradiction.}
	\bad{Je suppose $\sqrt3\notin\Rr$ et je cherche une contradiction.}
	\bad{Je suppose que $\sqrt3$ n'existe pas et je cherche une contradiction.}
\end{answers}
\begin{explanations}
Pour montrer que $\sqrt{3} \notin \Qq$ par l'absurde, on commence par supposer $\sqrt3\in\Qq$, donc on écrit $\sqrt{3}=\frac pq$ avec $p$ et $q$ entiers, puis on cherche une contradiction.
\end{explanations}
\end{question}
}
